% N3 WEN_COMPLEMENT --- research article of the FRONTEIRA_COMPACTA line.
% Gate: N0(c) (signature statement) + N0' (adversarial stress test) -- both green.
% Rev. 05jul26 (pos desk-reject CQG-116553): M5 classification theorem + decision
% table integrated (Sec. IV); M1c dichotomy upgrades CSG from measured to theorem;
% N6 quantum-robustness note in scope; reframed perspective -> research article.
% Claim discipline (charter): theorem about the CLASS (Axioms 1-2), never "Nature
% obeys"; every measured input cited to the campaign that measured it; open flanks
% declared in the final section.  Compiles with pdflatex+bibtex (revtex4-2).
\documentclass[aps,prd,reprint,nofootinbib,superscriptaddress,floatfix]{revtex4-2}
\usepackage{amsmath,amssymb}
\usepackage{bm}
\usepackage{booktabs}
\usepackage{hyperref}

\begin{document}

\title{A classification theorem for emergent $U(1)$ on causal substrates:\\
string-net admissibility coincides with broken Lorentz invariance}

\author{Miqueias Alves Mendes}
\email{mikalvesm@gmail.com}
\affiliation{Independent Researcher, Ibiapina, Cear\'a, Brazil}
\date{\today}

\begin{abstract}
Two research programmes derive gauge structure rather than postulate it. In condensed
matter, string-net condensation produces emergent photons as collective modes of loop
degrees of freedom on a lattice; sixty years of emergence phenomenology stand behind it.
In quantum gravity, causal set theory replaces spacetime by a Lorentz-invariant discrete
order; its kinematics recovers dimension, curvature and d'Alembertians---but no causal
substrate has ever produced an emergent photon. We show these two facts have a single
explanation, and that it is a theorem about signature rather than a technical gap.
String-net condensation needs three structures: finite valence, closed loops of finite
dimension, and an oriented spacelike 2-cell to carry magnetic flux. A Lorentz-invariant
causal substrate refuses all three, each by the non-compactness of the boost: every
invariant rule on a Poisson sprinkling has infinite expected valence (a Campbell--Mecke
theorem on the hyperboloid $H^{d-1}$); the measured escape routes fail a binary
alternative (finite valence \emph{or} finite-dimensional loops, never both); and the
causal order carries time with sign but space without sign and without structure, so the
magnetic 2-cell has no carrier---a prediction verified directly (interval `diamonds' are
$100\%$ electric). These pieces assemble into a classification theorem. Define a
substrate to be \emph{string-net admissible} (SNA) if its Hasse diagram has finite
valence, two-dimensional amenable ball growth, and a percolating plaquette density;
then \emph{a causal substrate is SNA if and only if it is non-invariant}: every
Lorentz-invariant or covariant-random class fails SNA by a proven mechanism, while the
deterministic crystal and the foliated lattice realize it---and a unified decision
table, one measured representative per class, confirms the boundary with no
counterexample. The exact boundary between the two worlds is the \emph{definiteness of
the metric signature}: separation orbits are compact spheres for $O(d)$ and infinite-volume
quadrics for $O(p,q)$, so divergence is \emph{chosen} on the Euclidean side and
\emph{forced} on the Lorentzian side. Wick rotation is precisely the operation that
crosses the boundary, and an order cannot cross it and remain an order. The statement
explains a negative: why condensed-matter emergence does not transport to causal
substrates---and it tells any future emergent-photon claim exactly which lemma it must
break.
\end{abstract}

\maketitle

\section{Two emergence programmes that never meet}

Emergence of gauge structure is no longer a speculation. In condensed matter it is a
theorem-grade phenomenology: local bosonic lattice models whose low-energy sectors contain
deconfined $U(1)$ gauge fields, with the photon as a collective mode of condensed
string-nets~\cite{LevinWen2005,LevinWenColloquium,WenBook,MotrunichSenthil2002}, and with
the robustness of the emergent gauge symmetry understood since
Foerster--Nielsen--Ninomiya~\cite{FNN1980}. In quantum gravity, the causal set programme
replaces the spacetime continuum by a locally finite partial
order~\cite{BLMS1987,Surya2019}; sprinkled at Poisson density it is the unique
Lorentz-invariant discretization~\cite{BHS2009}, and its kinematics demonstrably encodes
dimension, proper time, curvature, and horizon areas~\cite{Surya2019,Barton2019}.

The two programmes have an obvious meeting point---derive light on the causal
substrate---and a striking non-meeting record: no construction on a Lorentz-invariant
causal set has produced an emergent photon, ours included. Our own programme measured
this negative from several independent directions before understanding it (Sec.~III).
Here we show the non-meeting is structural: we state a classification theorem that
separates the two worlds (Sec.~IV), verify it on a unified decision table with one
measured representative per invariance class, and locate the boundary exactly: it is
the \emph{signature of the metric}, not the discreteness, not the dimension, and not a
lack of ingenuity.

Explaining a negative is rare and useful. The statement below explains simultaneously why
string-net emergence \emph{works} in condensed matter and why it \emph{cannot transport}
to Lorentz-invariant causal substrates---one theorem, two sides---and it converts years
of failed constructions into corollaries.

\section{What emergent light needs}

Stripped to its skeleton, string-net condensation~\cite{LevinWen2005,LevinWenColloquium}
requires of its substrate:

\begin{enumerate}
\item \textbf{Finite valence.} The degrees of freedom live on the links of a graph whose
vertices have finitely many neighbours; the Hamiltonian is a sum of local terms. Finite
valence is what makes ``local'' meaningful and Hilbert spaces finite-dimensional per
site.
\item \textbf{Loops of finite dimension.} The gauge-invariant configurations are closed
strings; the deconfined phase is a condensate of arbitrarily large loops built from
\emph{finite-dimensional} local loop moves (plaquettes). Without a finite plaquette
structure there is nothing to condense and no flux to quantize. (In $2{+}1$ dimensions
compact $U(1)$ still confines by monopole proliferation~\cite{Polyakov1977}; the
deconfined photon needs the loop structure \emph{and} enough dimensions---the
requirement is necessary, not sufficient.)
\item \textbf{An oriented spacelike 2-cell.} The magnetic field is a flux: $B$ is
defined by the holonomy around an oriented spacelike plaquette. A phase with photons
needs both electric (timelike) and magnetic (spacelike) 2-cells, with orientation.
\end{enumerate}

Condensed matter provides all three for free, because its substrate is a fixed spatial
lattice: a graph with finite valence, plaquettes by construction, and an orientation
inherited from the embedding in $\mathbb{R}^3$. The question is what a Lorentz-invariant
causal substrate can offer in their place.

\section{What the causal substrate refuses, layer by layer}

Fix the class: a Poisson sprinkling of $\mathbb{M}^d$ with its induced order
$(\mathcal{C},\prec)$, and any \emph{intrinsic} rule---a measurable functional of order
and counting only, no embedding data. This is the standard causal set kinematics; Poisson
is not optional but forced by invariance~\cite{BHS2009}. Three layers, in decreasing
generality; each kills one string-net requirement.

\subsection{Layer 1: infinite valence is forced (theorem)}

Any invariant pairwise quantity of a causally related pair is a function of the proper
time $\Delta\tau$ alone. The locus of ``neighbours at fixed $\Delta\tau$'' is the orbit
of the Lorentz group, the hyperboloid $H^{d-1}\!=\mathrm{SO}^+(d\!-\!1,1)/\mathrm{SO}(d\!-\!1)$,
of \emph{infinite} invariant volume. By the Campbell--Mecke formula, the expected degree
of any nontrivial invariant connection rule $q(\Delta\tau)$ is
\begin{equation}
\langle z\rangle \;=\; \rho\,\mathrm{Vol}(H^{d-1})\!\int_0^\infty\! \Delta\tau^{\,d-1}
q(\Delta\tau)\,d\Delta\tau \;=\; \infty ,
\label{eq:layer1}
\end{equation}
for every $q$ that is not almost-everywhere zero. The argument closes non-pairwise rules
as well (via the Palm calculus, up to the covering relation), and rules based on interval
cardinality $N(i,j)$ are $\Delta\tau$ in disguise
($\mathbb{E}[N|\Delta\tau]\propto\Delta\tau^d$ is a bijection). This is the familiar
non-locality of causal set links~\cite{Surya2019}, sharpened into an impossibility:
\emph{requirement 1 (finite valence) is unattainable for any invariant rule on the
sprinkling.} Not hard to attain---unattainable.

\subsection{Layer 2: the measured escape routes and the binary alternative}

Layer 1 has stated hypotheses (Poisson measure; rules up to the covering relation), and
we spent the programme's experimental effort trying to escape through them. The record
is a binary alternative, measured across seven independent substrate campaigns: a
substrate evades the \emph{boost} divergence or the \emph{loop} requirement, never both.

\begin{itemize}
\item \emph{Invariant repulsion} (Mat\'ern-II thinning in the invariant interval
$s^2$): $\langle z\rangle$ still diverges for every exclusion parameter---the
Lorentz-invariant exclusion volume is a non-compact band along the cone. Loops track
the Poisson control.
\item \emph{Invariant long-range percolation} $p(\Delta\tau)=\min\{1,(\Delta\tau/
\Delta\tau_0)^{-\sigma}\}$: for \emph{all} $\sigma$ scanned, $\langle z\rangle$ diverges
\emph{and} 4-cycles stay below the random control---the rule decays in the wrong
variable, because the invariant $\Delta\tau$ cannot cut the boost direction. The only
family that fails both barriers at once, everywhere in parameter space.
\item \emph{Non-Poisson growth} (classical sequential growth): the degree saturates
(escapes the boost) but the transitive reduction has provably zero clustering and
sub-mean-field 4-cycles---tree-like, no finite loop structure to condense. This exit
has since been closed as a theorem, not just a measurement: the covariant family is
geometrically dichotomous, hyperbolic or one-dimensional block, never two-dimensional
amenable (Sec.~IV).
\item \emph{Simplicial/CDT-like kinematics}: the only family whose 1-skeleton carries
saturating loops (4-cycles $\approx 0.145$, transitivity $\approx 0.30$ in the
two-dimensional flipped ensemble)---but its finite valence is the Euler identity of a
\emph{foliated} triangulation. Its three-dimensional and stacked four-dimensional
versions were measured mean-field (clustering decays as $N^{-0.33}$); and the foliation
itself is the point (next item).
\item \emph{Explicit foliation} (anisotropic, Ho\v{r}ava-type): the one substrate that
passed both barriers---by construction, because a preferred foliation \emph{breaks
Lorentz invariance}. The criticality it sustains was measured to be that of a stacked
three-dimensional Heisenberg lattice (exponent distance to the known lattice class
$0.89\sigma$): a known spatial lattice in disguise, i.e., the condensed-matter side of
the boundary, not an escape through it.
\end{itemize}

The alternative is exhaustive in everything measured so far: \emph{finite valence or
finite-dimensional loops---an invariant causal substrate never offers both.}
Requirement 2 falls.

\subsection{Layer 3: the magnetic 2-cell has no carrier}

Requirement 3 fails for a reason deeper than valence, and it was measured before it was
understood. The causal order carries \emph{time with sign} (the relation is directed;
interval counting estimates $\Delta\tau$ with its orientation) and \emph{space without
sign and without structure}: a spacelike pair contributes one bit (incomparability), and
every spatial datum beyond it is a statistical reconstruction from third
parties---unsigned Gram data, determining configurations at best up to the \emph{full}
orthogonal group $O(d{-}1)$, reflections included. Antichains, the discrete carriers of
``space,'' are order-theoretically amorphous: their only intrinsic datum is cardinality.

The electric field survives this: $E$ is flux through a \emph{timelike} 2-cell, and the
causal interval (``diamond'') is both structured (interval cardinality, longest chains)
and oriented (by the order itself). The magnetic field does not: $B$ is flux through an
oriented \emph{spacelike} 2-cell, which fails twice---no structure (amorphous
antichains) and no orientation (unsigned Gram data). The direct measurement agrees: on
sprinkled Minkowski, minimal interval 2-cells are $100\%$
electric~\cite{MatterGravityPRD}; background curvature opens a magnetic fraction only as
$\mathrm{frac}_B\propto H^2$---itself now derived, since the count is even under
time-reversal duality and must be even in $H$---with the flat-space value an exact zero.

\section{The classification theorem: the string-net boundary coincides with the
invariance boundary}

The three refusals of Sec.~III, together with the constructions that survive them,
assemble into a single decidable statement. Call a locally finite poset
\emph{string-net admissible} (SNA) if its Hasse diagram offers the three structures of
Sec.~II: \textbf{(SNA-1)} finite valence; \textbf{(SNA-2)} polynomial ball growth of
dimension ${\approx}\,2$---an amenable, lattice-like skeleton rather than a
hyperbolic expander; \textbf{(SNA-3)} a percolating $\Theta(N)$ density of minimal
plaquettes. The predicate is intrinsic---order and counting only, no embedding
data---and each clause is decidable by measurement on any substrate ensemble.

\medskip
\noindent\textbf{Theorem (classification).}
\emph{(Exclusion: SNA $\Rightarrow$ non-invariant.)} No Lorentz-invariant or
covariant-random substrate is SNA. Invariant rules on the Poisson sprinkling fail
SNA-1 by Eq.~(\ref{eq:layer1}); exchangeable random orders possess no locally finite
Hasse diagram at all; embedded box-orders fail SNA-1 logarithmically; and covariant
sequential growth obeys a geometric dichotomy---sparse regimes are hyperbolic
(tree-like: SNA-2 and SNA-3 fail), dense regimes are one-dimensional block chains
punctuated by posts (SNA-2 fails)---so it is never two-dimensional amenable.
\emph{(Realizability: SNA is inhabited in the non-invariant class.)} The deterministic
crystal $\mathbb{Z}^2$ and the foliated square lattice are SNA; both lie outside every
invariance class (no crystal is Lorentz-invariant, and square-generated percolating
loops force a grading, hence a preferred foliation).

\medskip
\noindent\textbf{Corollary (the decision boundary).} \emph{A causal substrate is
string-net admissible if and only if it is non-invariant.} An emergent photon on a
causal substrate requires breaking invariance; there is no invariant route.

\medskip
Table~\ref{tab:sna} is the theorem made empirical: one representative of every named
invariance class, measured on the same axis with the same estimators. SNA is true
exactly on the two non-invariant lines; each invariant line fails by the mechanism its
theorem names; no counterexample appeared. The converse direction is existence, not
universality: not every non-invariant order is SNA (dust is not)---the content of the
corollary is that the invariant side is \emph{empty} and the non-invariant side is
\emph{inhabited}.

\begin{table*}[t]
\caption{\label{tab:sna}The unified decision table. Columns: mean Hasse valence
$\langle z\rangle$ ($\uparrow\!N$ marks divergence with system size); post density;
ball-growth exponential rate and polynomial degree (amenability window: rate $<0.35$
with degree $\in[1.6,2.4]$, measured at the largest resolvable radius); minimal
4-cycles per node (plaquette percolation). SNA-1/2/3 as defined in the text.
Exchangeable orders are excluded by theorem (no locally finite Hasse diagram), hence
carry no measured row. SNA is true exactly on the two non-invariant lines.}
\begin{ruledtabular}
\begin{tabular}{lcccccccccc}
Class & invariant? & $\langle z\rangle$ & posts & rate & poly & $C_4/N$ &
SNA-1 & SNA-2 & SNA-3 & SNA \\
\colrule
Poisson $\mathbb{M}^2$ & yes (Lorentz) & $11\;(\uparrow\!N)$ & $0$ & $0.69$ & $2.70$
 & $29.6$ & $\times$ & $\times$ & $\checkmark$ & \textbf{no} \\
box-order 2D & yes (embedded) & $11\;(\uparrow\!N)$ & $0$ & $0.82$ & $2.81$ & $29$
 & $\times$ & $\times$ & $\checkmark$ & \textbf{no} \\
CSG sparse & covariant & $3.9$ & $0$ & $0.66$ & $3.32$ & $0.02$
 & $\checkmark$ & $\times$ & $\times$ & \textbf{no} \\
CSG dense & covariant & $2.7$ & $0.35$ & $0.14$ & $0.99$ & $0.56$
 & $\checkmark$ & $\times$ & $\checkmark$ & \textbf{no} \\
exchangeable & yes & --- & --- & --- & --- & ---
 & $\times$ & $\times$ & $\times$ & \textbf{no} \\
crystal $\mathbb{Z}^2$ & \textbf{no} & $9.6$ & $0$ & $0.27$ & $1.93$ & $18.4$
 & $\checkmark$ & $\checkmark$ & $\checkmark$ & \textbf{yes} \\
foliated lattice & \textbf{no} & $4.0$ & $0$ & $0.25$ & $1.76$ & $1.95$
 & $\checkmark$ & $\checkmark$ & $\checkmark$ & \textbf{yes} \\
\end{tabular}
\end{ruledtabular}
\end{table*}

\section{The boundary is the signature}

Each layer above traces to one fact: the boost is non-compact. That fact has an exact
converse, and together they locate the boundary between the two emergence worlds.

Let a Poisson process in $(\mathbb{R}^d, g)$ be invariant under the isometries of a
metric of signature $(p,q)$, with any nontrivial invariant connection rule. The orbit of
neighbours at fixed invariant separation is the quadric $\{x\cdot x=c\}$:

\begin{itemize}
\item \textbf{Definite signature} ($pq=0$; isometry group $O(d)$, compact): the orbit is
the sphere $S^{d-1}$, of finite volume. $\langle z\rangle<\infty$ for every integrable
profile; finite valence, geometric random graphs, genuine criticality, string-nets,
photons. Divergence is possible but \emph{chosen} (a long-tailed profile).
\item \textbf{Indefinite signature} ($p,q\geq 1$; $O(p,q)$ non-compact): \emph{every}
separation quadric has infinite invariant volume, and Eq.~(\ref{eq:layer1}) applies
verbatim: $\langle z\rangle=\infty$ for \emph{every} invariant rule. Divergence is
\emph{forced}.
\end{itemize}

The boundary is not ``Lorentzian versus Euclidean'' but \emph{definite versus
indefinite}; the Lorentzian case $q=1$ is simply the unique indefinite signature whose
cone is convex enough to define an order. The asymmetry is
\emph{forced-versus-chosen}, which is why no Euclidean counterexample can exist: on the
definite side one can always choose a compactly supported invariant profile.

This also says exactly what Wick rotation is, seen from the substrate: the operation
that replaces $O(d-1,1)$ by $O(d)$---that \emph{compactifies the group}. Euclidean
lattice gauge theory and condensed matter live on the definite side from the start;
that is why their emergence machinery works. A causal order cannot Wick-rotate and
remain an order: the order \emph{is} the Lorentzian datum. There is no rotation that
carries a causal set to the side where string-nets condense.

We note a precedent inside standard physics for reading the boundary this way: in
Wigner's classification the discrete quantum numbers come from the compact little group,
the continuous labels from the non-compact factors, and the mass shell that carries the
continuous momentum label is the \emph{same} hyperboloid $H^{d-1}$ whose infinite volume
drives Eq.~(\ref{eq:layer1}). One geometric object, two manifestations of one
non-compactness.

\section{Consequences}

\textbf{For emergent-photon programmes.} Any claim of an emergent $U(1)$ on a causal
substrate must now specify which hypothesis it breaks, because inside the class the
route is closed at three independent layers. The options are exhaustive and each is a
declared research programme, not a loophole: (i) \emph{break Lorentz invariance}
(preferred foliation)---measured to reproduce known spatial-lattice universality, i.e.,
to re-enter condensed matter rather than extend it; (ii) \emph{leave the Poisson
measure} for genuinely dynamical geometry---for covariant sequential growth this exit
is now closed by the geometric dichotomy of Sec.~IV; what remains open is genuinely
non-Markovian or quantum-dynamical geometry, where our strongest constraint is only
measured (mean-field behaviour in every CDT-like ensemble tested), not proven;
(iii) \emph{import the gauge sector}---the
honest default of the causal set programme itself, which couples matter fields to the
order rather than deriving them.

\textbf{For horizon thermodynamics.} The same boundary has an entropic face, measured
recently in this programme: on sprinkled Minkowski the \emph{geometric} horizon count
(Barton \emph{et al.} molecules~\cite{Barton2019,DouSorkin2003}) obeys a clean area law
in $d=2,3,4$---while the mutual information of an ordered \emph{matter} field across
the same Rindler corner grows super-area ($\sim L^{3.4}$ against $L^2$ at accessible
sizes), because the field lives on links whose corner-crossing number itself grows
super-area. Geometry, built from the order, keeps the area law; matter, built on the
non-local links, loses it. The non-compactness is not a technical nuisance of link
counting; it is one boundary with connectivity, criticality, and entropy faces.

\textbf{Falsifiability.} The statement is attackable at named joints: exhibit a
Lorentz-invariant rule with finite expected valence (breaks Layer 1 and with it the
Campbell--Mecke argument); or an invariant substrate with saturating
finite-dimensional loops (breaks the measured binary); or an intrinsic construction of
an oriented spacelike 2-cell from order and counting (breaks Layer 3 and the unsigned-Gram
lemma behind it). Any one of these would be more interesting than an emergent photon.

\section{Scope, and what remains open}

The claim is a theorem about a class---Poisson causal sets with intrinsic rules---not
about Nature; it is kinematic, in the same sense in which Wigner's classification is
kinematic. Its known gaps are part of the statement. Layer 1 is rigorous for pairwise
rules and for the non-pairwise closure via Palm calculus, up to the covering relation.
The loop side of the binary (Layer 2), a measured pattern across seven substrate families
when this line began, has since been upgraded to a trichotomy of theorems: no exchangeable
random order has a locally finite Hasse diagram; growth orders with posts confine every
cycle to a finite block; and square-generated percolating loops force a grading, hence a
foliation incompatible with boost invariance. The one combinatorial escape the theorems
left open---percolating \emph{rank-frustrated} (pentagonal) loops---has now been
constructed explicitly, but only as rigid crystals lying outside every invariance class,
so it confirms rather than threatens the boundary: the obstruction is invariance, not the
causal order as such. Covariant sequential growth is closed by the geometric dichotomy
of Sec.~IV; what escapes the classification is genuinely non-Markovian growth, the one
named open flank. Two robustness checks bound the statement from outside the classical
kinematics: the boundary is insensitive to replacing classical statistical weights by
interfering amplitudes over the same ensembles (the SNA verdict by valence class
reappears unchanged in a first quantum crossing, indicating the obstruction lives in
the order, not in the weight), and dynamical geometry in the CDT universality class
remains constrained empirically (mean-field in every ensemble tested). Within the
class, the conclusion is closed: the absence of an emergent
photon on Lorentz-invariant causal sets is not an open problem awaiting a cleverer
construction. It is the signature, doing to gauge structure exactly what it does to
valence, to criticality, and---we now know---to the entropy of matter across a horizon.

\begin{acknowledgments}
The measured inputs cited as companion work are reported in
Refs.~\cite{MatterGravityPRD,SU3PRD}, with full pre-registrations, death criteria and
negative results in the programme's public record.
\end{acknowledgments}

\bibliography{wen_complement}

\end{document}
